% ===========================================================================
%  Concentration of Measure  --  proof notes organized by engine
%  woven on LOOM.   compile:  xelatex main  (twice)  or  latexmk -xelatex main
%  How to fill: each section ships thesis / paradigm / cheat-sheet / theorem
%  statements + proof skeletons; replace the \TODO{...} and the skeleton stubs.
% ===========================================================================
\documentclass[cjk]{loom}   % [cjk] only so the pen name renders; all prose is English

% --- \block \trigger \TODO + booktabs/tabularx/L now come from loom.cls v0.2 ---
\newcommand{\R}{\mathbb R}
\newcommand{\E}{\mathbb E}
\newcommand{\Var}{\operatorname{Var}}
\newcommand{\Ent}{\operatorname{Ent}}
\DeclareMathOperator{\tr}{tr}

\runningthread{concentration of measure}

\begin{document}

\loomcover
  {Concentration of Measure}
  {Proof notes, organized by engine}
  {北极甜虾 (剑心犹在！)}
  {\today}

% ===========================================================================
\section*{Overview · the eight engines}
\addcontentsline{toc}{section}{Overview · the eight engines}
\warmth{0}

\begin{strand}
In one line: a function of many \emph{weakly dependent} variables is nearly
constant. Every concentration inequality is one \keyword{engine} for proving
this; they differ only in the structure they consume: independence, bounded
differences, convexity, or curvature and isoperimetry. Learn the engine before
the formula.
\end{strand}

\trigger{You want a small-probability bound on $f(X_1,\dots,X_n)$ deviating far from its mean/median.}
Ask yourself one question: \emph{which structure do I have?} Then reach for the engine below.

\block{Cheat-sheet: which engine to use}
\noindent\begin{tabularx}{\linewidth}{@{}l l L@{}}
\toprule
{\color{indigo}Engine} & {\color{indigo}Structure it eats} & \multicolumn{1}{l}{\color{indigo}Main tools / section} \\
\midrule
Cramér–Chernoff        & independence + MGF            & Laplace transform, optimize $\lambda$, rate function \;(§\ref{sec:chernoff}) \\
Sub-Gaussian / sub-exp.& \emph{shape} of the tail      & $\psi$-Orlicz norms, Bernstein \;(§\ref{sec:subgaussian}) \\
Bounded diff.\ / martingale & change one coordinate    & Azuma, McDiarmid \;(§\ref{sec:bdd}) \\
Variance / entropy     & tensorization                 & Efron–Stein, Herbst, log-Sobolev \;(§\ref{sec:entropy}) \\
Geometry / isoperimetry& metric \emph{and} measure     & sphere/Gaussian isoperimetry, Borell–TIS \;(§\ref{sec:geometry}) \\
Talagrand convex dist. & convexity + product space     & convex-distance inequality \;(§\ref{sec:talagrand}) \\
Matrix concentration   & spectrum + trace exponential  & Lieb, matrix Bernstein \;(§\ref{sec:matrix}) \\
Suprema / chaining     & metric entropy                & Dudley, generic chaining \;(§\ref{sec:chaining}) \\
\bottomrule
\end{tabularx}

\begin{strand}
Reading guide: §\ref{sec:chernoff}–§\ref{sec:subgaussian} are the trunk (the
continuation of Appendix A); read them until fluent. §\ref{sec:bdd}–§\ref{sec:entropy}
push ``sums'' to ``any bounded-difference function''.
§\ref{sec:geometry}–§\ref{sec:talagrand} are the geometric heart that gives
``\emph{measure} concentration'' its name. §\ref{sec:matrix}–§\ref{sec:chaining}
are two exits toward ML/TCS.
\end{strand}

\bigskip
{\small\tableofcontents}

\clearpage
% ===========================================================================
% ===========================================================================
\section{The Cramér–Chernoff engine}\label{sec:chernoff}
\warmth{0}\quad\whisper{The trunk. Every later engine just works on one step of this pipeline.}

\begin{strand}
In one line: to bound the probability that a random quantity strays far from its
mean by something \keyword{exponentially small}, exponentiate the tail through the
Laplace transform (the MGF) and pick a good $\lambda$ to push the right-hand side
down. The whole subject is variations on this single move.
\end{strand}

\block{The paradigm}
\trigger{You want an exponentially small bound on $\Pr[X\ge a]$, with the MGF of $X$ under control.}
For every $\lambda>0$,
\[
  \Pr[X\ge a]=\Pr\!\big[e^{\lambda X}\ge e^{\lambda a}\big]\le \E[e^{\lambda X}]\,e^{-\lambda a}
  \qquad(\text{Markov on }e^{\lambda X}),
\]
then take the infimum over $\lambda$. Writing the cumulant generating function
$\psi_X(\lambda)=\log\E[e^{\lambda X}]$,\warp{mgf}
\[
  \log\Pr[X\ge a]\le -\psi_X^*(a),\qquad
  \psi_X^*(a)=\sup_{\lambda>0}\big(\lambda a-\psi_X(\lambda)\big),
\]
where $\psi_X^*$ is the \keyword{Legendre transform} (the rate function) of $\psi_X$.

\begin{strand}
Three-line heuristic: (1) independence $\Rightarrow$ MGFs \emph{multiply}, which is the
\emph{only} reason a ``sum'' gets an exponential bound; (2) \emph{center} first,
$X\leftarrow X-\E X$, for the cleanest argument; (3) bound each $\E[e^{\lambda X_i}]$ by one
\emph{termwise convex} estimate. Every theorem in the subject just swaps that estimate
and that $\lambda$.
\end{strand}

\block{Cheat-sheet: the three opening moves}
\noindent\begin{tabularx}{\linewidth}{@{}l l L@{}}
\toprule
{\color{indigo}Tool} & {\color{indigo}Setting} & \multicolumn{1}{l}{\color{indigo}Tail bound / when to use}\\
\midrule
Markov           & $Y\ge 0$                     & $\Pr[Y>\alpha\E Y]<1/\alpha$; opening move, first-order info only\\
Chebyshev        & $\Var X<\infty$              & $\Pr[|X-\E X|>t]\le \Var X/t^2$; second-order info\\
Chernoff (master)& MGF exists                   & $\Pr[X\ge a]\le e^{-\psi_X^*(a)}$; want an exponential tail\\
\bottomrule
\end{tabularx}

\block{Main results \& proof skeletons}

\begin{theorem}[Chernoff master bound]\label{thm:chernoff-master}
If the MGF of $X$ is finite near $0$, then $\displaystyle \Pr[X\ge a]\le e^{-\psi_X^*(a)}$.
\end{theorem}
\begin{proof}
Skeleton: \emph{(i)} Markov on $e^{\lambda X}$; \emph{(ii)} take logs to get $\log\Pr[X\ge a]\le \psi_X(\lambda)-\lambda a$;
\emph{(iii)} take the infimum over $\lambda>0$. \TODO{add convexity of $\psi_X$, the attaining $\lambda$, and the equality case}
\end{proof}

\begin{lemma}[Additivity: independence makes $\psi$ add]\label{lem:mgf-add}
$X_1,\dots,X_n$ independent $\Rightarrow \psi_{\sum_i X_i}=\sum_i \psi_{X_i}$.
\end{lemma}
\begin{proof}\TODO{one line: $\E[e^{\lambda\sum X_i}]=\prod\E[e^{\lambda X_i}]$, take logs}\end{proof}

\begin{example}[The sub-Gaussian seed: a symmetric $\pm1$ sum]\label{ex:rademacher}
Let $\varepsilon_i$ be independent uniform $\pm1$ and $S_n=\sum_{i\le n}\varepsilon_i$. From
$\cosh\lambda\le e^{\lambda^2/2}$ we get $\psi_{S_n}(\lambda)\le \lambda^2 n/2$, and $\lambda=a/n$ gives
\[
  \Pr[S_n\ge a]\le e^{-a^2/2n},\qquad \Pr[|S_n|\ge a]\le 2e^{-a^2/2n}.
\]
\TODO{prove $\cosh\lambda\le e^{\lambda^2/2}$ by termwise Taylor comparison; this is the seed inequality of the whole subject}
\end{example}

\begin{strand}
Why it goes this way: $\cosh\lambda\le e^{\lambda^2/2}$ is the \emph{seed inequality}. Replace it by a
different termwise estimate (Hoeffding's lemma $\to$ bounded variables, Bennett/Bernstein
$\to$ variance-aware) and you grow §\ref{sec:subgaussian}–§\ref{sec:entropy}. Swap the estimate, swap $\lambda$, that is all.
\end{strand}

\block{Tightness: Chernoff gets the large-deviation rate right}
\begin{theorem}[Cramér, matching lower bound]\label{thm:cramer}
Under mild conditions $\displaystyle \lim_{n\to\infty}\tfrac1n\log\Pr\!\big[\tfrac1n\textstyle\sum_i X_i\ge a\big]=-\psi_{X_1}^*(a)$.
\end{theorem}
\begin{proof}
Skeleton: the upper bound is Thm~\ref{thm:chernoff-master}+Lem~\ref{lem:mgf-add}; the lower bound uses
\emph{exponential tilting} (change of measure) to make $a$ the mean under the new measure, then CLT.
\TODO{construct the tilted measure $d\mathbb{Q}\propto e^{\lambda x}\,d\mathbb{P}$}
\end{proof}

\loose{The three deviation regimes (keep this picture): small deviations $\to$ normal; large deviations $\lambda_n=o(\sigma_n)\to \log\Pr\sim-\lambda_n^2/2$; very large deviations $\lambda_n=\Omega(\sigma_n)\to$ normal fails, Poisson takes over. Cf.\ Appendix A.}%
\pick{mgf} This thread is upgraded in §\ref{sec:subgaussian} to ``sub-Gaussian $\equiv$ MGF dominated by $e^{\sigma^2\lambda^2/2}$''.

% ===========================================================================
\section{Sub-Gaussian and sub-exponential}\label{sec:subgaussian}
\warmth{0}\quad\whisper{Turn ``what does the tail look like'' into a norm. Once classified, tail bounds come for free.}

\begin{strand}
In one line: rather than recompute MGFs each time, classify random variables by the
\emph{shape of their tail}. Sub-Gaussian = tail no heavier than a Gaussian's; sub-exponential
= no heavier than an exponential's. Each family gets an Orlicz norm, and concentration
becomes arithmetic on those norms.
\end{strand}

\block{The paradigm}
\trigger{You want to handle a whole class of tails at once, and have the tail of sums and products follow \emph{mechanically}.}
Four equivalent characterizations (sub-Gaussian, write $\|X\|_{\psi_2}$): tail
$\Pr[|X|>t]\le 2e^{-t^2/K^2}$ $\Leftrightarrow$ moments $\|X\|_p\lesssim K\sqrt p$ $\Leftrightarrow$
MGF $\E e^{\lambda X}\le e^{K^2\lambda^2/2}$ (after centering) $\Leftrightarrow$ $\E e^{X^2/K^2}\le 2$.

\begin{strand}
Three-line heuristic: (1) \keyword{a centered bounded variable is sub-Gaussian} (Hoeffding's
lemma), so the entire ``bounded-difference'' family lands in range here; (2) sums of
sub-Gaussians stay sub-Gaussian and the proxy variances \emph{add}:
$\|\sum X_i\|_{\psi_2}^2\le\sum\|X_i\|_{\psi_2}^2$; (3) a product of two sub-Gaussians is sub-\emph{exponential},
which is exactly why covariance / quadratic-form estimates drop to the sub-exponential family.
\end{strand}

\block{Cheat-sheet: the two tail families}
\noindent\begin{tabularx}{\linewidth}{@{}l l L@{}}
\toprule
{\color{indigo}Theorem} & {\color{indigo}Setting} & \multicolumn{1}{l}{\color{indigo}Tail bound / when to use}\\
\midrule
Hoeffding lemma   & $\E X=0$, $X\in[a,b]$            & $\E e^{\lambda X}\le e^{\lambda^2(b-a)^2/8}$; bridge bounded $\to$ sub-Gaussian\\
Hoeffding ineq.\  & $X_i\in[a_i,b_i]$ independent    & $\Pr[|S-\E S|\ge t]\le 2e^{-2t^2/\sum(b_i-a_i)^2}$\\
General Hoeffding & $X_i$ sub-Gaussian               & $\Pr[|\sum a_iX_i|\ge t]\le 2e^{-ct^2/\|a\|_2^2K^2}$\\
Bernstein         & $X_i$ sub-exp., $\E X_i=0$       & $\Pr[|S|\ge t]\le 2\exp\!\big[-c\min(\tfrac{t^2}{\sum\nu_i^2},\tfrac{t}{\max b_i})\big]$\\
Bennett           & $|X_i|\le b$, variance $\sigma^2$& variance-aware; Gaussian for small $t$, Poisson for large $t$\\
\bottomrule
\end{tabularx}

\block{Main results \& proof skeletons}

\begin{definition}[Sub-Gaussian norm]\label{def:subg}
$\|X\|_{\psi_2}=\inf\{K>0:\ \E\exp(X^2/K^2)\le 2\}$. Call $X$ \keyword{sub-Gaussian} if this is finite.
(Sub-exponential $\|X\|_{\psi_1}$ replaces $X^2$ by $|X|$.)
\end{definition}

\begin{proposition}[Four equivalences]\label{prop:subg-equiv}
For centered $X$ the four characterizations above are pairwise equivalent, with constants differing only by absolute factors.
\end{proposition}
\begin{proof}
Skeleton: tail $\Rightarrow$ moments (integrate the tail) $\Rightarrow$ MGF (sum the series $\sum K^p p^{p/2}\lambda^p/p!$)
$\Rightarrow$ tail (Chernoff, §\ref{sec:chernoff}) $\Rightarrow$ Orlicz. \TODO{track the constants through the four steps; draw it as a cycle}
\end{proof}

\begin{lemma}[Hoeffding's lemma]\label{lem:hoeffding}
$\E X=0$, $X\in[a,b]$ a.s.\ $\Rightarrow \E e^{\lambda X}\le e^{\lambda^2(b-a)^2/8}$.
\end{lemma}
\begin{proof}
Skeleton: set $\psi(\lambda)=\log\E e^{\lambda X}$; compute $\psi(0)=\psi'(0)=0$ and
$\psi''(\lambda)=\Var_{\lambda}(X)\le(b-a)^2/4$ (the variance under the tilted measure is at most a quarter of the range squared); Taylor remainder.
\TODO{prove $\psi''\le(b-a)^2/4$ via the variance-range estimate}
\end{proof}

\begin{theorem}[Hoeffding's inequality]\label{thm:hoeffding}
$X_i\in[a_i,b_i]$ independent, $S=\sum X_i$. Then $\Pr[|S-\E S|\ge t]\le 2\exp\!\big(-\tfrac{2t^2}{\sum_i(b_i-a_i)^2}\big)$.
\end{theorem}
\begin{proof}\TODO{feed Lem~\ref{lem:hoeffding} into the master bound of §\ref{sec:chernoff} + Lem~\ref{lem:mgf-add}, optimize $\lambda$}\end{proof}

\begin{theorem}[Bernstein's inequality]\label{thm:bernstein}
$X_i$ independent, centered, sub-exponential. Then the tail is Gaussian for small $t$ and exponential for large $t$ (see the cheat-sheet row).
\end{theorem}
\begin{proof}
Skeleton: the single-term sub-exponential MGF bound $\E e^{\lambda X_i}\le e^{\nu_i^2\lambda^2/2}$ holds only for
$|\lambda|<1/b_i$, so the choice of $\lambda$ is constrained by $1/b$, producing the two-regime $\min$.
\TODO{do the constrained $\lambda$-optimization (the boundary between the two regimes)}
\end{proof}

\begin{strand}
Why it goes this way: a Gaussian tail comes from ``MGF dominated everywhere by $e^{c\lambda^2}$'';
an exponential tail comes from ``dominated only for $|\lambda|<1/b$''. The constrained $\lambda$-optimization
automatically yields Bernstein's $\min(t^2/\nu,\,t/b)$ split, which is the whole secret of being \emph{variance-aware}.
\end{strand}

\begin{example}[Johnson–Lindenstrauss, preview]\label{ex:jl}
The squared length of a Gaussian random projection $\frac1{\sqrt m}Gx$ is a $\chi^2$-type sum (sub-exponential),
and Bernstein gives $\Pr[|\,\|\Pi x\|^2-\|x\|^2\,|>\epsilon\|x\|^2]\le 2e^{-cm\epsilon^2}$.
\TODO{add the union bound giving $m=O(\epsilon^{-2}\log N)$}
\end{example}

\loose{Check the constant: does Hoeffding's $2t^2/\sum(b-a)^2$ agree with Appendix A's A.1.18 (range $\le1$ gives $e^{-2a^2/n}$)?}%
\pick{mgf}\quad
\warp{subgaussian} The sub-Gaussian thread is reused throughout §\ref{sec:bdd} (McDiarmid) and §\ref{sec:matrix} (matrix Bernstein) via \pick{subgaussian}.

% ===========================================================================
\section{Bounded differences and martingales}\label{sec:bdd}
\warmth{0}\quad\whisper{Upgrade ``sum of independents'' to ``any function that changes little when one coordinate moves''.}

\begin{strand}
In one line: it need not be a sum. As long as $f(x_1,\dots,x_n)$ changes by a bounded amount
when a single coordinate is altered, reveal its randomness \emph{one coordinate at a time}
(the Doob martingale); each step has bounded difference, feed it to Hoeffding's lemma from
§\ref{sec:subgaussian}, and concentration appears.
\end{strand}

\block{The paradigm}
\trigger{$f$ is not a sum, but has \keyword{bounded differences}: altering coordinate $i$ changes it by $|f-f'|\le c_i$.}
Build the \keyword{Doob martingale} $M_k=\E[f\mid X_1,\dots,X_k]$, so $M_0=\E f$, $M_n=f$, with
increments $|M_k-M_{k-1}|\le c_k$. \pick{subgaussian}Apply Hoeffding's lemma per step and sum.

\begin{strand}
Three-line heuristic: (1) \keyword{any} function is its own Doob martingale (revealing one coordinate = one step);
(2) concentration depends only on \emph{how far each step can jump}, not on how complex $f$ is;
(3) the martingale constant ($\sum c_i^2$) is looser than the entropy method (§\ref{sec:entropy}) by a factor,
but the proof is shortest and the assumptions weakest.
\end{strand}

\block{Cheat-sheet: the bounded-difference family}
\noindent\begin{tabularx}{\linewidth}{@{}l l L@{}}
\toprule
{\color{indigo}Theorem} & {\color{indigo}Setting} & \multicolumn{1}{l}{\color{indigo}Tail bound / when to use}\\
\midrule
Azuma–Hoeffding & increments $|M_k-M_{k-1}|\le c_k$ & $\Pr[|M_n-M_0|\ge t]\le 2e^{-t^2/2\sum c_k^2}$\\
McDiarmid       & $f$ bounded diff.\ $c_i$, $X_i$ indep.\ & $\Pr[|f-\E f|\ge t]\le 2e^{-2t^2/\sum c_i^2}$; workhorse of ML generalization\\
Self-bounding   & $f$ self-bounding              & one-sided Bernstein type, variance-aware (cf.\ §\ref{sec:entropy})\\
Freedman        & martingale + predictable var.\ $V$ & variance-aware martingale tail, tighter when $V$ small\\
\bottomrule
\end{tabularx}

\block{Main results \& proof skeletons}

\begin{theorem}[Azuma–Hoeffding]\label{thm:azuma}
$(M_k)$ a martingale, $|M_k-M_{k-1}|\le c_k$ a.s. Then $\Pr[M_n-M_0\ge t]\le \exp\!\big(-t^2/2\textstyle\sum_k c_k^2\big)$.
\end{theorem}
\begin{proof}
Skeleton: for the increments $D_k=M_k-M_{k-1}$ we have $\E[D_k\mid\mathcal F_{k-1}]=0$ and $|D_k|\le c_k$,
so the conditional Hoeffding lemma gives $\E[e^{\lambda D_k}\mid\mathcal F_{k-1}]\le e^{\lambda^2 c_k^2/2}$;
peel off layer by layer via the tower property to get $\E e^{\lambda(M_n-M_0)}\le e^{\lambda^2\sum c_k^2/2}$; then Chernoff.
\TODO{fill the conditional Hoeffding lemma and the tower-property unrolling}
\end{proof}

\begin{theorem}[McDiarmid's bounded-differences inequality]\label{thm:mcdiarmid}
$X_1,\dots,X_n$ independent, $f$ with bounded differences $c_i$. Then $\Pr[|f-\E f|\ge t]\le 2\exp\!\big(-2t^2/\textstyle\sum_i c_i^2\big)$.
\end{theorem}
\begin{proof}\TODO{feed the Doob martingale $M_k=\E[f\mid X_{1:k}]$ into Thm~\ref{thm:azuma}; independence keeps each increment's range $\le c_k$ (sup minus inf)}\end{proof}

\begin{example}[Concentration of the generalization gap]\label{ex:generalization}
$f=\sup_{h\in\mathcal H}\big(\E_{\mathcal D}\ell_h-\widehat{\E}\,\ell_h\big)$ has bounded differences $\le 1/n$ in each sample,
so McDiarmid gives $\Pr[f-\E f\ge t]\le e^{-2nt^2}$. Thus the generalization gap concentrates uniformly around its
expectation ($=$ the Rademacher complexity). \TODO{verify the bounded-difference constant $1/n$}
\end{example}

\begin{strand}
Why it goes this way: McDiarmid splits ``uniform convergence'' into two things, \emph{concentration}
(this section, $f$ hugs $\E f$) and \emph{the size of the expectation} ($\E f=$ Rademacher complexity, estimated
by the chaining of §\ref{sec:chaining}). That is the source of the standard two-step recipe in statistical learning theory.
\end{strand}

\loose{Bounded differences can be relaxed: the asymmetric / variance-aware refinement (BLM Ch.\,6) replaces $\sum c_i^2$ by the true variance, but needs the entropy method. Room to improve here.}%
\warp{tensorize} The truly tight bound comes from the \emph{tensorization} of §\ref{sec:entropy}; this section is its ``martingale approximation''. \pick{tensorize}

% ===========================================================================
\section{The variance and entropy methods}\label{sec:entropy}
\warmth{0}\quad\whisper{Tensorization: variance and entropy are subadditive over product measures, so high-dim = one-dim subadditivity + a single-coordinate estimate.}

\begin{strand}
In one line: both variance (Efron–Stein) and entropy (log-Sobolev / Herbst) are \emph{subadditive}
over product measures. \keyword{Tensorize} a high-dimensional concentration into $n$ single-coordinate
contributions, each needing only a one-dimensional estimate. This route gives the tightest,
\emph{variance-aware, dimension-free} bounds, and is the analytic engine behind the geometry of
§\ref{sec:geometry} and Talagrand (§\ref{sec:talagrand}).
\end{strand}

\block{The paradigm}
\trigger{You want the tightest \emph{variance-aware} bound, or to prove a function sub-Gaussian without a ``sum / martingale'' structure.}
\keyword{Herbst's argument}: if the entropy obeys $\Ent[e^{\lambda f}]\le \tfrac{\lambda^2\sigma^2}{2}\,\E[e^{\lambda f}]$
for all $\lambda$, then $f-\E f$ is $\sigma$-sub-Gaussian. Combine with the \keyword{tensorization} of entropy
$\Ent[g]\le \sum_i \E\,\Ent_i[g]$ to reduce the problem to a single coordinate.

\begin{strand}
Three-line heuristic: (1) \keyword{Efron–Stein} is the tensorization of variance, giving the \emph{free}
(no i.i.d.\ needed) bound $\Var f\le \tfrac12\sum_i\E(f-f^{(i)})^2$; (2) replace $\Var$ by $\Ent$ and ``square''
by ``exponential'': tensorization + a one-dimensional log-Sobolev $=$ exponential concentration (Herbst);
(3) the Gaussian measure satisfies a \emph{dimension-free} log-Sobolev inequality (Gross), the algebraic source
of the Gaussian concentration in §\ref{sec:geometry}.
\end{strand}

\block{Cheat-sheet: variance and entropy}
\noindent\begin{tabularx}{\linewidth}{@{}l l L@{}}
\toprule
{\color{indigo}Theorem} & {\color{indigo}Setting} & \multicolumn{1}{l}{\color{indigo}Gives / when to use}\\
\midrule
Efron–Stein     & product measure, $f^{(i)}=$ resample coord.\ $i$ & $\Var f\le\tfrac12\sum_i\E(f-f^{(i)})^2$; the ``bounded difference'' for variance\\
Bounded diff.\ via ES & bounded diff.\ $c_i$               & $\Var f\le\tfrac14\sum c_i^2$; variance from the martingale family\\
Herbst argument & $\Ent[e^{\lambda f}]\le\tfrac{\lambda^2\sigma^2}2\E e^{\lambda f}$ & $\Rightarrow f$ is $\sigma$-sub-Gaussian\\
Gaussian LSI    & $\gamma=N(0,I_n)$               & $\Ent[g^2]\le 2\E\|\nabla g\|^2$; dimension-free\\
Modified LSI    & bounded / Bernoulli coords      & tensorize + single-coord LSI $\Rightarrow$ concentration\\
\bottomrule
\end{tabularx}

\block{Main results \& proof skeletons}

\begin{theorem}[Efron–Stein inequality]\label{thm:efron-stein}
$X_1,\dots,X_n$ independent, $X_i'$ an i.i.d.\ copy, $f^{(i)}$ replacing coordinate $i$ by $X_i'$. Then
$\displaystyle \Var f\le \tfrac12\sum_{i=1}^n\E\big[(f-f^{(i)})^2\big]$.
\end{theorem}
\begin{proof}
Skeleton: use the orthogonal martingale-difference decomposition $\Var f=\sum_i\E[\Delta_i^2]$ with
$\Delta_i=\E[f\mid X_{1:i}]-\E[f\mid X_{1:i-1}]$; then Jensen to replace $\Delta_i$ by the symmetrized $(f-f^{(i)})$.
\TODO{fill the orthogonal decomposition + symmetrization, two steps}
\end{proof}

\begin{theorem}[Herbst's argument]\label{thm:herbst}
If $\Ent[e^{\lambda f}]\le\tfrac{\lambda^2\sigma^2}2\E[e^{\lambda f}]$ for all $\lambda$, then
$\E[e^{\lambda(f-\E f)}]\le e^{\lambda^2\sigma^2/2}$, so $f$ is $\sigma$-sub-Gaussian.
\end{theorem}
\begin{proof}
Skeleton: let $H(\lambda)=\tfrac1\lambda\log\E e^{\lambda f}$; the entropy inequality is exactly the rewriting
$H'(\lambda)\le\sigma^2/2$; integrate to get $H(\lambda)-H(0^+)\le\lambda\sigma^2/2$, with $H(0^+)=\E f$.
\TODO{fill the differential identity $\Ent[e^{\lambda f}]=\lambda^2\E e^{\lambda f}\,H'(\lambda)$}
\end{proof}

\begin{theorem}[Gaussian log-Sobolev, Gross]\label{thm:gaussian-lsi}
For $\gamma=N(0,I_n)$ and smooth $g$, $\Ent_\gamma[g^2]\le 2\,\E_\gamma\|\nabla g\|^2$.
\end{theorem}
\begin{proof}\TODO{fill: the one-dimensional LSI (Ornstein–Uhlenbeck semigroup, or CLT limit from two-point Bernoulli) + tensorization to $n$ dimensions}\end{proof}

\begin{strand}
Why it goes this way: Gaussian LSI $+$ Herbst $\Rightarrow$ ``Gaussian concentration $e^{-t^2/2L^2}$ for an
$L$-Lipschitz function, dimension-free''. This is the algebraic path that \emph{computes} the geometric
statement of §\ref{sec:geometry}: curvature (the LSI constant) turns directly into the tail rate.
\end{strand}

\loose{Four faces: concentration $\Leftrightarrow$ isoperimetry $\Leftrightarrow$ log-Sobolev $\Leftrightarrow$ transportation (Marton). Herbst is the ``LSI $\Rightarrow$ concentration'' face; the other three are to be filled, see §\ref{sec:geometry} and §9 open threads.}%
\pick{tensorize}\quad Tensorization reappears in §\ref{sec:talagrand} in Marton's transportation form.

% ===========================================================================
\section{Geometry and isoperimetry}\label{sec:geometry}
\warmth{0}\quad\whisper{Where ``concentration of measure'' gets its name: on a metric measure space, a Lipschitz function is almost constant.}

\begin{strand}
In one line: a viewpoint with no MGF at all. On a metric measure space, every $1$-Lipschitz function
\keyword{hugs its median}; this is equivalent to an \keyword{isoperimetric} phenomenon, a set of measure
$\ge\tfrac12$, blown up by $t$, swallows $1-e^{-ct^2}$ of the measure. High-dimensional spheres and Gaussian
space behave exactly so; this is the geometric root of ``a random vector's length is almost determined'' and
``a random projection preserves distance''.
\end{strand}

\block{The paradigm}
\trigger{$f$ is an $L$-Lipschitz function on a space with \keyword{concentration function} $\alpha(t)$.}
Then $\Pr[|f-\operatorname{Med}f|>t]\le 2\,\alpha(t/L)$. Three pillar spaces:
\[
  \begin{aligned}
    \text{sphere } S^{n-1}:&\ \alpha(t)\le e^{-(n-1)t^2/2};\\
    \text{Gaussian } \gamma_n:&\ \alpha(t)\le e^{-t^2/2}\quad(\text{\keyword{dimension-free}});\\
    \text{Hamming cube}:&\ \alpha(t)\le e^{-2t^2/n}.
  \end{aligned}
\]

\begin{strand}
Three-line heuristic: (1) concentration $\equiv$ isoperimetry: the tail bound is the complement of ``the
$t$-blow-up of a half-measure set''; (2) median and mean are interchangeable (differ by $O(L)$, same order);
(3) the Gaussian case being \emph{dimension-free} is the key miracle, it makes infinite-/high-dimensional
processes (§\ref{sec:chaining}) possible.
\end{strand}

\block{Cheat-sheet: geometric concentration}
\noindent\begin{tabularx}{\linewidth}{@{}l l L@{}}
\toprule
{\color{indigo}Theorem} & {\color{indigo}Space} & \multicolumn{1}{l}{\color{indigo}Statement / when to use}\\
\midrule
Lévy's lemma    & sphere $S^{n-1}$          & $1$-Lip $f$: $\Pr[|f-\operatorname{Med}f|>t]\le 2e^{-(n-1)t^2/2}$\\
Gaussian conc.\ & $\gamma_n=N(0,I_n)$       & $L$-Lip $f$: $\Pr[|f-\E f|>t]\le 2e^{-t^2/2L^2}$, \emph{dimension-free}\\
Borell–TIS      & Gaussian process         & $\sup_t X_t$ concentrates around its supremum, variance $=\sup_t\Var X_t$\\
Gaussian isoper.& $\gamma_n$                & half-spaces are extremal; exact blow-up in $\Phi^{-1}$ form\\
\bottomrule
\end{tabularx}

\block{Main results \& proof skeletons}

\begin{theorem}[Lévy's lemma / spherical concentration]\label{thm:levy}
$f:S^{n-1}\to\R$ is $1$-Lipschitz. Then $\Pr[|f-\operatorname{Med}f|>t]\le 2e^{-(n-1)t^2/2}$.
\end{theorem}
\begin{proof}
Skeleton: spherical \keyword{isoperimetry} (Lévy–Schmidt: geodesic balls are extremal) $\Rightarrow$ for a
half-measure set $A$, the $t$-neighborhood satisfies $\sigma(A_t)\ge 1-e^{-(n-1)t^2/2}$; take $A=\{f\le\operatorname{Med}f\}$,
and Lipschitz pushes $\{f>\operatorname{Med}f+t\}$ outside $A_t$. \TODO{fill the standard isoperimetry $\Rightarrow$ concentration translation; cite spherical isoperimetry (do not reprove)}
\end{proof}

\begin{theorem}[Gaussian concentration / Borell–TIS]\label{thm:gaussian-conc}
$f:\R^n\to\R$ is $L$-Lipschitz, $X\sim N(0,I_n)$. Then $\Pr[|f(X)-\E f|>t]\le 2e^{-t^2/2L^2}$, with constants independent of $n$.
\end{theorem}
\begin{proof}
Either path: \emph{(a)} Gaussian log-Sobolev (Thm~\ref{thm:gaussian-lsi}) $+$ Herbst (Thm~\ref{thm:herbst}), with smooth
approximation of $f$; or \emph{(b)} Gaussian isoperimetry (half-space extremality) gives the blow-up directly.
\TODO{take (a): fill the Lipschitz-gradient equivalence $\|\nabla f\|\le L$ a.e.\ + smoothing}
\end{proof}

\begin{example}[A random vector's length is almost determined]\label{ex:norm-conc}
$x\mapsto\|x\|$ is $1$-Lipschitz $\Rightarrow$ for $X\sim N(0,I_n)$, $\|X\|$ concentrates near $\sqrt n$ with
fluctuation $O(1)$ (independent of $n$). This is the one-line proof of the ``thin-shell'' phenomenon.
\TODO{fill $\E\|X\|=\sqrt n(1-o(1))$ and substitute}
\end{example}

\begin{strand}
Why it goes this way: the MGF method (§\ref{sec:chernoff}) answers ``sums''; the geometric method answers
``\emph{any} Lipschitz function'', dimension-free in the Gaussian case, at the cost of needing a ``hard'' input
like isoperimetry / log-Sobolev. The two paths \emph{meet} at Gaussian concentration: the LSI of §\ref{sec:entropy}
is precisely the algebraic proof of Borell–TIS.
\end{strand}

\loose{Median vs mean: they differ by $O(L)$, but exact isoperimetry is most natural with the median. When must they be distinguished? (heavy tails, asymmetric $f$.)}%
\warp{borelltis} Borell–TIS (concentration of the supremum of a Gaussian process) is the ``concentration half'' of the chaining in §\ref{sec:chaining}, chaining controls the \emph{expectation}, this theorem the \emph{fluctuation}. \pick{borelltis}

% ===========================================================================
\section{Talagrand's convex distance}\label{sec:talagrand}
\warmth{0}\quad\whisper{A product space has no geometry, only convexity. Talagrand's convex distance adapts to ``how many bad coordinates''.}

\begin{strand}
In one line: on a pure \emph{product} probability space (no sphere-like geometry), Talagrand's
\keyword{convex distance} yields dimension-free, very strong concentration for convex Lipschitz functions.
Its edge over bounded differences (§\ref{sec:bdd}): it does not look at the ``largest single-coordinate jump''
but at ``how \emph{many} coordinates you must change to escape a set'', adapting to sparse dependence.
\end{strand}

\block{The paradigm}
\trigger{$f$ on a product space is \keyword{convex} Lipschitz, or escaping an event $A$ requires changing many coordinates.}
Define the \keyword{convex distance} $d_T(x,A)=\sup_{\|\alpha\|_2\le1}\ \inf_{y\in A}\sum_{i:x_i\ne y_i}\alpha_i$ (weighted Hamming).
Talagrand: $\E\big[e^{d_T(X,A)^2/4}\big]\le 1/\Pr[A]$, hence
\[
  \Pr[A]\cdot\Pr[d_T(X,A)\ge t]\le e^{-t^2/4}.
\]

\begin{strand}
Three-line heuristic: (1) $d_T$ is a min-max of \emph{weighted} Hamming, automatically finding the ``cheapest escape direction'';
(2) once $\Pr[A]\ge\tfrac12$, the convex distance from $A$ concentrates sub-Gaussianly;
(3) for a \emph{convex} $L$-Lipschitz $f$, take $A=\{f\le\operatorname{Med}f\}$ to get $\Pr[|f-\operatorname{Med}f|>t]\le 4e^{-t^2/4L^2}$, dimension-free.
\end{strand}

\block{Cheat-sheet: convex distance and its corollaries}
\noindent\begin{tabularx}{\linewidth}{@{}l l L@{}}
\toprule
{\color{indigo}Result} & {\color{indigo}Setting} & \multicolumn{1}{l}{\color{indigo}Statement / when to use}\\
\midrule
Convex-distance ineq.\ & product space, any $A$    & $\Pr[A]\Pr[d_T\ge t]\le e^{-t^2/4}$; the master switch\\
Convex Lipschitz    & $f$ convex, $L$-Lip on $[0,1]^n$ & $\Pr[|f-\operatorname{Med}f|>t]\le 4e^{-t^2/4L^2}$\\
LIS / combinatorial opt.\ & $f=$ longest incr.\ subseq., etc.\ & $\sqrt{\,}$-scale fluctuation, beats bounded diff.'s $n$-scale\\
Largest singular value & random matrix          & $\sigma_{\max}$ concentrates around its median\\
\bottomrule
\end{tabularx}

\block{Main results \& proof skeletons}

\begin{theorem}[Talagrand's convex-distance inequality]\label{thm:talagrand}
$X=(X_1,\dots,X_n)$ a product measure, $A$ measurable. Then $\E\big[e^{d_T(X,A)^2/4}\big]\le 1/\Pr[A]$.
\end{theorem}
\begin{proof}
Skeleton (induction on $n$, the hardest step in the book): \emph{(i)} write $d_T$ via convex duality (the distance to
the convex hull $U_A(x)=\operatorname{conv}\{\mathbf 1[x_i\ne y_i]:y\in A\}$); \emph{(ii)} condition on the last coordinate
and apply the inductive hypothesis to the two cases ``in $A$ / projection in $A$''; \emph{(iii)} a one-dimensional
optimization $\inf_{0\le r\le1}r^{\dots}\cdots$ merges the two cases. \TODO{fill the three-step induction; the one-dim optimization lemma $\inf_r e^{(1-r)^2/4}(\cdot)^{r}$ is the crux}
\end{proof}

\begin{corollary}[Concentration of convex Lipschitz functions]\label{cor:convex-lip}
$f:[0,1]^n\to\R$ convex and $L$-Lipschitz (Euclidean), $X$ a product measure. Then
$\Pr[|f-\operatorname{Med}f|>t]\le 4e^{-t^2/4L^2}$.
\end{corollary}
\begin{proof}\TODO{use convexity to control $f(x)-\operatorname{Med}f$ by $L\cdot d_T(x,\{f\le\operatorname{Med}f\})$ (convex $\Rightarrow$ supporting hyperplane), then apply Thm~\ref{thm:talagrand}}\end{proof}

\begin{example}[Longest increasing subsequence (LIS)]\label{ex:lis}
For a random permutation, $L_n\sim 2\sqrt n$; changing one value moves it by at most $1$, but ``certifying $L_n\ge k$ needs
only $k$ positions'', so the convex distance gives fluctuation $O(n^{1/4})$, far beating bounded differences' $O(\sqrt n)$.
\TODO{fill the certifiable / self-certifying structure $\Rightarrow d_T$ bound}
\end{example}

\begin{strand}
Why it goes this way: bounded differences only knows ``each coordinate has $c_i$''; Talagrand knows ``how \emph{many}
coordinates must conspire to change the conclusion''. When dependence is sparse (few coordinates determine $f$), convex
distance replaces $\sqrt{\sum c_i^2}$ by $\sqrt{\#\text{critical coords}}$, the source of the $n\to\sqrt n\to n^{1/4}$
chain of improvements in combinatorial optimization.
\end{strand}

\loose{Another route: Marton's transportation-cost ($W_1\le\sqrt{2\,\mathrm{KL}}$) also gives convex-distance concentration, tying it back to the entropy tensorization of §\ref{sec:entropy}. Worth listing side by side.}%
\pick{tensorize}\quad Convex distance can also be seen as the scalar precursor of the spectral concentration of §\ref{sec:matrix}.

% ===========================================================================
\section{Matrix concentration}\label{sec:matrix}
\warmth{0}\quad\whisper{Lift scalar Chernoff/Bernstein to the spectral norm of a random matrix. The obstacle is $e^{A+B}\ne e^Ae^B$.}

\begin{strand}
In one line: to bound the largest eigenvalue of a sum of independent random matrices, copying §\ref{sec:chernoff}
stalls because \emph{matrix exponentials do not multiply}. The rescue is \keyword{Lieb's concavity} $+$ the trace
exponential: use $\E\,\tr\exp(\cdot)$ as a surrogate; it is \emph{subadditive}, so the scalar proof transfers almost
verbatim, with the dimension entering only at the end as $\log d$.
\end{strand}

\block{The paradigm}
\trigger{$\sum_i X_i$ is a sum of independent random Hermitian matrices and you want a tail for $\|\sum_i X_i\|$ (spectral norm).}
\keyword{Matrix Laplace transform}: for $\theta>0$,
\[
  \Pr\big[\lambda_{\max}(\textstyle\sum_i X_i)\ge t\big]\le e^{-\theta t}\,\E\,\tr\exp\!\Big(\textstyle\sum_i \log\E\,e^{\theta X_i}\Big),
\]
where the $\tr\exp(\sum\cdots)$ comes from Lieb's theorem ($A\mapsto\tr\exp(H+\log A)$ is concave) $+$ Jensen.

\begin{strand}
Three-line heuristic: (1) turn ``$\lambda_{\max}\ge t$'' into $\tr e^{\theta\sum X_i}\ge e^{\theta t}$ then Markov;
(2) \keyword{Lieb} lets us push $\E$ inside $\tr\exp$ and makes the matrix cumulants \emph{subadditive} (replacing the scalar
``MGFs multiply''); (3) the dimension $d$ enters only through $\tr\le d\cdot\lambda_{\max}$, so it is nearly dimension-free
at \emph{low intrinsic dimension}.
\end{strand}

\block{Cheat-sheet: matrix tail bounds}
\noindent\begin{tabularx}{\linewidth}{@{}l l L@{}}
\toprule
{\color{indigo}Theorem} & {\color{indigo}Setting} & \multicolumn{1}{l}{\color{indigo}Tail bound / when to use}\\
\midrule
Matrix Laplace  & independent Hermitian $X_i$    & the display above; mother of all matrix tails\\
Matrix Chernoff & $0\preceq X_i\preceq R\,I$      & $\lambda_{\min/\max}(\sum X_i)$ around $\mu_{\min/\max}$; PSD sums\\
Matrix Bernstein& $\E X_i=0,\|X_i\|\le L$, $\|\sum\E X_i^2\|=\nu$ & $\Pr[\|\sum X_i\|\ge t]\le 2d\,e^{-t^2/2(\nu+Lt/3)}$\\
Intrinsic-dim.\ & as above, $\mathrm{intdim}=\nu/\|\cdot\|$ & $d\to$ intrinsic dim, drops the ambient-dimension $\log$\\
\bottomrule
\end{tabularx}

\block{Main results \& proof skeletons}

\begin{theorem}[Matrix Laplace transform method, Tropp]\label{thm:matrix-laplace}
For independent random Hermitian $X_i$ and $\theta>0$, the matrix Laplace inequality above holds.
\end{theorem}
\begin{proof}
Skeleton: \emph{(i)} $\Pr[\lambda_{\max}\ge t]\le e^{-\theta t}\E\,\tr e^{\theta\sum X_i}$ (Markov $+$ $\lambda_{\max}\le\tr$ on PSD);
\emph{(ii)} \keyword{Lieb's theorem} $\Rightarrow A\mapsto \tr\exp(H+\log A)$ concave $\Rightarrow$ Jensen pushes $\E$ inside;
\emph{(iii)} peel off term by term to get the subadditive cumulant $\tr\exp(\sum_i\log\E e^{\theta X_i})$.
\TODO{fill the Lieb $\Rightarrow$ subadditivity induction; cite Lieb (do not reprove)}
\end{proof}

\begin{theorem}[Matrix Bernstein]\label{thm:matrix-bernstein}
$X_i$ independent, $\E X_i=0$, $\|X_i\|\le L$, variance proxy $\nu=\|\sum_i\E X_i^2\|$. Then
$\Pr[\|\sum_i X_i\|\ge t]\le 2d\exp\!\big(\tfrac{-t^2/2}{\nu+Lt/3}\big)$.
\end{theorem}
\begin{proof}\TODO{feed the single-term matrix cgf bound $\log\E e^{\theta X_i}\preceq \frac{\theta^2/2}{1-L\theta/3}\E X_i^2$ into Thm~\ref{thm:matrix-laplace}, optimize $\theta$}\end{proof}

\begin{example}[Covariance estimation]\label{ex:covariance}
For the sample covariance $\widehat\Sigma$ of $n$ independent samples, write $\widehat\Sigma-\Sigma=\frac1n\sum_i X_i$ (centered outer products);
matrix Bernstein gives $\|\widehat\Sigma-\Sigma\|\lesssim\sqrt{\tfrac{\|\Sigma\|\,r\log d}{n}}$ ($r=$ intrinsic dimension).
\TODO{fill the $L,\nu$ estimates of the outer products and the sample-size threshold}
\end{example}

\begin{strand}
Why it goes this way: the entire structure of scalar Bernstein (the variance term $\nu$ $+$ the uniform bound $L$,
the two-regime $\min$) survives untouched, with only an extra $d$ (or intrinsic-dim) prefactor, coming from turning
``one tail'' into ``the tail of $d$ eigenvalues''. This machine replaces the ``union bound over an $\epsilon$-net''
common in ML proofs, and is both tighter and less work.
\end{strand}

\loose{Martingale version: matrix Freedman (Tropp 2011) extends this to matrix martingales, pairing with §\ref{sec:bdd}. Needed for the spectral concentration of online / adaptive sampling.}%
\warp{matrixbern} This thread feeds directly into sketching and leverage-score sampling. \pick{subgaussian}

% ===========================================================================
\section{Suprema and chaining}\label{sec:chaining}
\warmth{0}\quad\whisper{The exit. When the target is $\sup_t X_t$ (an empirical process), single-point Chernoff + a union bound is too loose.}

\begin{strand}
In one line: the first seven sections control \emph{one} random quantity. But ML/statistics really wants
$\sup_{t\in T}X_t$ (empirical processes, the sup in a generalization bound). \keyword{Chaining} uses multi-scale nets
to control the supremum by an integral of metric entropy, reducing ``infinitely many tails'' to ``the geometry of
covering numbers''. This is the bridge to empirical processes.
\end{strand}

\block{The paradigm}
\trigger{You want $\E\sup_{t\in T}X_t$, where $\{X_t\}$ is a (sub-)Gaussian process with natural metric $d(s,t)=\|X_s-X_t\|_{\psi_2}$.}
\keyword{Chaining}: along a sequence of refining $\varepsilon$-nets $T_0\subset T_1\subset\cdots$, telescope $X_t-X_{t_0}$ and sum scale by scale. \keyword{Dudley}:
\[
  \E\sup_{t\in T}X_t\ \le\ C\int_0^\infty\sqrt{\log N(T,d,\varepsilon)}\ d\varepsilon.
\]

\begin{strand}
Three-line heuristic: (1) a union bound $=$ a single scale, chaining $=$ \emph{all scales} summed, hence tighter;
(2) the currency is \keyword{metric entropy} $\log N(T,d,\varepsilon)$, turning a probability problem into covering geometry;
(3) Dudley can be loose; the \keyword{generic chaining} quantity $\gamma_2(T,d)$ is the \emph{tight} one (Talagrand's majorizing measures).
\end{strand}

\block{Cheat-sheet: from union bound to majorizing measures}
\noindent\begin{tabularx}{\linewidth}{@{}l l L@{}}
\toprule
{\color{indigo}Tool} & {\color{indigo}Gives} & \multicolumn{1}{l}{\color{indigo}When to use / tight?}\\
\midrule
Union bound     & $\sqrt{2\log|T|}$ (finite $T$)  & opening move; too loose for large or continuous $T$\\
Dudley integral & $\int\sqrt{\log N(\varepsilon)}\,d\varepsilon$ & standard upper bound for continuous index; often enough\\
Sudakov lower bd.\ & $\sup_\varepsilon\varepsilon\sqrt{\log N(\varepsilon)}$ & pairs with Dudley to pin the order\\
Generic chaining& $\gamma_2(T,d)$               & matching upper \emph{and} lower (Gaussian processes, Talagrand)\\
\bottomrule
\end{tabularx}

\block{Main results \& proof skeletons}

\begin{theorem}[Dudley's inequality]\label{thm:dudley}
$\{X_t\}_{t\in T}$ a centered process, sub-Gaussian w.r.t.\ $d$. Then
$\E\sup_{t}X_t\le C\int_0^{\infty}\sqrt{\log N(T,d,\varepsilon)}\,d\varepsilon$.
\end{theorem}
\begin{proof}
Skeleton: take nets $T_k$ at scale $\varepsilon_k=2^{-k}$ with nearest-point maps $\pi_k(t)$; telescope
$X_t-X_{\pi_0(t)}=\sum_{k\ge1}(X_{\pi_k(t)}-X_{\pi_{k-1}(t)})$; each link has $|T_k||T_{k-1}|$ sub-Gaussian increments,
combine the per-scale union bound with the increment scale $\lesssim 2^{-k}$ and sum.
\TODO{fill the per-scale union bound $\sqrt{\log|T_k|}$ and the geometric-series sum $\Rightarrow$ the integral}
\end{proof}

\begin{theorem}[Generic chaining / majorizing measures]\label{thm:generic-chaining}
For a Gaussian process, $\E\sup_t X_t\asymp\gamma_2(T,d)$, where
$\gamma_2(T,d)=\inf_{(T_k)}\sup_t\sum_{k\ge0}2^{k/2}d(t,T_k)$ with $|T_k|\le 2^{2^k}$.
\end{theorem}
\begin{proof}\TODO{upper bound: weighted chaining (weights adapt to $t$, beating Dudley's uniform scales); lower bound: Talagrand's partition tree / iterated Sudakov. Cite the majorizing-measures theorem (do not reprove)}\end{proof}

\begin{example}[Operator norm via $\varepsilon$-net + chaining]\label{ex:operator-norm}
$\|A\|=\sup_{u,v\in S^{n-1}}u^\top Av$: discretize on a $1/2$-net of the sphere ($\le 9^n$ points) + each point sub-Gaussian
(§\ref{sec:subgaussian}) $\Rightarrow$ $\E\|A\|\lesssim\sqrt n$. Chaining ``automates'' the net argument to a continuous index.
\TODO{fill the standard $1/2$-net lemma lifting $\sup$ to $2\max_{\text{net}}$}
\end{example}

\begin{strand}
Why it goes this way: the Borell–TIS of §\ref{sec:geometry} (\pick{borelltis}) controls the \emph{fluctuation} of $\sup_t X_t$,
this section its \emph{expectation}, together they are the full concentration of an empirical process. This closes the loop
with the ``expectation $=$ Rademacher complexity'' half of the generalization bound in §\ref{sec:bdd}.
\end{strand}

\loose{This is the exit: metric entropy / empirical processes / uniform laws deserve \emph{their own notebook} (Dudley, van der Vaart–Wellner, Talagrand's \emph{Upper and Lower Bounds}). This section only goes as far as ``chaining + Borell–TIS'', as a bridge.}%
\warp{chaining} Chaining $\to$ Rademacher complexity $\to$ uniform laws of large numbers.

% ===========================================================================
\section*{Open threads · where this goes next}
\addcontentsline{toc}{section}{Open threads · where this goes next}
\warmth{0}

\begin{strand}
These notes are an \keyword{arsenal}, not a museum. The cross-section open questions below are the
seeds for the next pass: a unifying perspective, the missing lower bounds, and the exit toward
empirical processes.
\end{strand}

\block{Open threads (cross-section)}
\begin{itemize}[leftmargin=1.4em,itemsep=2pt]
  \item \keyword{Four faces}: concentration $\Leftrightarrow$ isoperimetry $\Leftrightarrow$ log-Sobolev
        $\Leftrightarrow$ transportation (Marton). Drawing the explicit commuting diagram of these equivalences
        is the ``main theorem'' this notebook most wants to add.
        \loose{Once this diagram is in, §\ref{sec:entropy}–§\ref{sec:talagrand} unify into one engine.}
  \item \keyword{When does dimension enter}? Gaussian concentration is dimension-free, matrix concentration carries
        $\log d$, chaining carries metric entropy. A ``sources of dimension dependence'' table separates ``truly
        dimension-free'' from ``hidden in the constants''.
  \item \keyword{Tightness}: is every upper bound paired with a lower bound? Cramér (§\ref{sec:chernoff}) and
        Sudakov (§\ref{sec:chaining}) are; the lower bounds for McDiarmid and matrix Bernstein (worst $f$ / worst spectrum) are to be filled.
  \item \keyword{The exit}: after §\ref{sec:chaining}, open a dedicated \emph{empirical-process} notebook (VC / Rademacher / uniform laws / Donsker).
\end{itemize}

\bigskip
\begin{strand}
Closing line: the whole subject is seven variations on one move, \emph{exponentiate the deviation, then crush it
with structure}. What the structure is decides which engine you reach for. Recognizing the structure beats memorizing the formula.
\end{strand}


\end{document}
